\scanpage{25}
\mspage{13}
\margin{24}

Il ne semble pas que la structure de graphe orienté cellularisé pondéré de $\Gamma^*(\underline\Sigma)$ permette de distinguer, parmi les \deleted{sommets} arêtes issues d'un sommet $D$ i.e. les sommets de $\operatorname{Pol}^*(D)$, celles et ceux qui correspondent à des triangles de ps. spécialisation (parmi les \deleted{\unsure{paires de sommets}} \inserted{sommets incidents} à des arêtes \deleted{\unsure{de poids}} de $\operatorname{Pol}^*(D)$ de poids nul, dont \deleted{\unsure{l'antipodique est un sommet}} l'un au moins soit sommet de générisation).

\margin{mais si ! cf ci-contre !}

Cette distinction peut se déduire si on dispose des applications de spécialisation \deleted{\unsure{et des}} $\operatorname{Pol}^*(D)\longrightarrow\operatorname{Pol}^*(D')$ : les sommets de ps. spécialisation sont distingués, parmi les précédents, comme \deleted{aux} les ``spécialisés'' des sommets de spécialisation véritables.

La prochaine question est donc si les applications de spécialisation
\[
\operatorname{Pol}^*(D)\longrightarrow\operatorname{Pol}^*(D')
\]
sont déductibles des structures précédentes. \uline{Mais bien sûr}, via l'\deleted{effet} isom correspondant sur les polygônes déployés, car il suffit de

\unsure{a)} savoir \unsure{qu'ils sont} de $=$ ordre (\unsure{calculable|découle})

b) de savoir \deleted{l'effet} du morphisme sur l'ens des deux orientations, ce qui est connu, et

c) de savoir l'image d'un sommet particulier \inserted{$(D,a)$} de $\widehat{\operatorname{Pol}}^*(D)$, qui est \unsure{connu} comme l'antipodique, dans $\widehat{\operatorname{Pol}}(D')$, du sommet $(D',a)$ !

En fait, chaque fois qu'on a une 2-carte \inserted{(avec bouts isolés i.e. arêtes libres)} dont les \deleted{\unsure{deux}} (vrais) sommets sont d'ordre $2$, et dont les arêtes et coins sont pondérés [avec pondérations des arêtes $\ge0$], \deleted{\unsure{et où}} et la pondération totale de chaque sommet \deleted{est} (somme des pondérations des \deleted{\unsure{inc}} \inserted{arcs} et coins incidents) est cte, on trouve des isom de transport entre les polygônes locaux déployés, qui forment donc un syst. local sur la surface sous-jacente à la carte !
